\documentclass[conference]{IEEEtran}
\usepackage{cite}
\usepackage{amsmath,amssymb,amsfonts}
\usepackage{algorithmic}
\usepackage{graphicx}
\usepackage{textcomp}
\usepackage{xcolor}

\begin{document}

\title{Empirical Analysis of Uniform and Log-Adaptive Bucket Sort on Right-Skewed Commercial Dataset Distributions}

\author{\IEEEauthorblockN{Makayla}
\IEEEauthorblockA{\textit{Department of Computer Science} \\
GitHub Repository: \texttt{vgsales-bucket-sort-research}}
}

\maketitle

\begin{abstract}
Comparison-based sorting algorithms possess a theoretical lower bound of $\mathcal{O}(n \log n)$. Distribution sorts such as Bucket Sort achieve $\mathcal{O}(n + k)$ average-case complexity under uniform distributions. However, real-world commercial datasets—exemplified by video game global sales—exhibit heavy right-skewness, leading to catastrophic clustering and $\mathcal{O}(n^2)$ worst-case performance in naive implementations. This paper evaluates standard uniform bucket sorting against a $\log(1+x)$ adaptive transformation variant using the Kaggle Video Game Sales dataset, analyzing empirical runtime and peak memory consumption across scaling input sizes.
\end{abstract}

\section{Introduction}
Sorting forms a foundational pillar in data processing pipelines. While comparison sorts like Timsort operate universally, distribution sorts exploit structural knowledge of input ranges. This work explores how bucket allocation strategies impact performance when processing power-law distributed numerical values.

\section{Methodology \& Algorithm Design}
We implemented two variations of Bucket Sort:
\begin{enumerate}
    \item \textbf{Standard Uniform Bucket Sort:} Maps continuous floating-point ranges linearly into $k$ buckets.
    \item \textbf{Log-Adaptive Bucket Sort:} Applies a logarithmic transformation $\log(1+x)$ prior to bucket mapping to mitigate extreme right-skew clustering.
\end{enumerate}

\section{Conclusion}
Empirical benchmarks confirm that log-adaptive bucket scaling effectively prevents worst-case runtime collapse on heavy-tailed distributions.

\end{document}
